\section{\classlabelshort Basic Setup and General Guide} 

\subsection{\classlabelshort Scenario Setup}  
%%%%% BEGIN USER CONTENT %%%%%
\guidance{
	\unclass Do not change this section. 
}
\classlabelshort Please reference the \gls{odessa} User and Analyst Manual for details on how to configure the \gls{odessa} \gls{gui}. This section is focused on the \modelname{} \gls{gui}.

\subsection{\classlabelshort Model Specific Setup}  
%\figHere{GUI Selection}{../images/GUIConfiguration1.png}
\guidance{
\unclass TIMELINE: This section should be completed by \gls{trr}, but any plans to modify the \gls{gui} should be described/pictured here no later than \gls{cdr} even if the work hasn't been completed. 

\unclass \gls{gui} image below should be replaced with your system. Use multiple images and incremental descriptions if necessary. Zoom in to make it easier to read. 

\unclass For \gls{odessa}, several models require input CSVs for waveforms, etc.  The set up of these files needs to be included in this section.

\unclass Specific text around the system options that can be specified in the GUI should be shown in \pic{GUI Options}. 

}
\doMyFigure{7in}{\unclasslabel}{../images/GUIConfiguration1}{\unclass GUI Selection}{GUI Selection}

%%%%% BEGIN USER CONTENT %%%%%
\unclass The \modelname{} is selected from the "asset selection" tab within the GUI shown in \pic{GUI Selection}. The system can be added by either double clicking on the \modelname, or it can be selected and then click "add to scenario". 

\doMyFigure{2.75in}{\unclasslabel}{../images/placeholder}{\unclass GUI System Options}{GUI Options}




%%Input Parameters
\subsection{\unclass  Model Parameters}  
\guidance{
	\unclass TIMELINE: This section should be filled in for \gls{pdr} and finalized for \gls{cdr}.
	
\unclass The purpose of this section is to provide a consolidated location for all of the input parameters used to make this a system-specific model, whether hard-coded or set by the user.

\unclass The Component column should show which component is most directly affected by the setting.

\unclass Refer to sources by reference number from the References Table.

\unclass Options for user-tunable column are no (it's hard-coded), \gls{gui}, or super user input. 

\unclass Also include settings that may need to be set to specific values such as the time step in the \gls{gui}. 

}

\unclass Parameters used or available for the model are listed in \tref{ParameterTable}. The user-tunable column has three options: No, \gls{gui} input, or super-user input (SUI).  

\begin{center}
	\begin{longtblr}
		[
		caption={\unclass Model Parameters},
		label={tab:ParameterTable},
		]{
			colspec={|Q[valign=m,l,1in]|Q[valign=m,l,1.5in]|Q[valign=m,l,1in]|Q[valign=m,c,0.75in]|Q[valign=m,c,0.625in]|Q[valign=m,c,0.7in]|},
			cell{1}{1}={c=6}{halign=l,bg=white,fg=black},
			cell{Z}{1}={c=6}{halign=r,bg=white,fg=black},
			row{2}={bg=TableBlue,fg=white,halign=c},
			hline{3}={.5pt},
			hline{Y}={.5pt},
			rowhead=2,
			rowfoot=1,
		}   \hline
		\unclasslabel \\\hline
		Variable & Description & Default and Classification & Component & User Tunable? & Source Number \\
		OperationalBW & Frequency range the system operates over & \unclass 100-200MHz & Receiver & No & 2 \\ \hline
		Frequency & Center frequency of the system & \unclass 110MHz & Receiver & \gls{gui} & 1 \\
		\unclasslabel \\\hline
	\end{longtblr}
\end{center}


